% Generated from locale/tr/content/many-valued-logic/syntax-and-semantics/connectives.tex; do not hand-edit.


\olfileid[tr]{mvl}{syn}{con}

\olsection{Diller ve Bağlaçlar}

Klasik önerme mantığı ve başka birçok mantık, \emph{önerme sabitleri} ile
\emph{bağlaçlardan} oluşan sabit bir dağarcık kullanır. Örneğin
aşağıdakileri ilkel kabul ederiz:
\begin{enumerate}
\tagitem{prvFalse}{!!{falsity} için önerme sabiti $\lfalse$.}{}
\tagitem{prvTrue}{!!{truth} için önerme sabiti $\ltrue$.}{}
\item Mantıksal bağlaçlar:
  \startycommalist
  \iftag{prvNot}{\ycomma $\lnot$ (değilleme)}{}%
  \iftag{prvAnd}{\ycomma $\land$ (ve bağlacı)}{}%
  \iftag{prvOr}{\ycomma $\lor$ (veya bağlacı)}{}%
  \iftag{prvIf}{\ycomma $\lif$ (!!{conditional})}{}%
  \iftag{prvIff}{\ycomma $\liff$ (!!{biconditional})}{}%
\end{enumerate}
\iftag{defNot,defOr,defAnd,defIf,defIff,defTrue,defFalse,defEx,defAll}{%
Yukarıdaki ilkel bağlaçlara ek olarak, kısaltma olarak tanımlanan şu
simgeleri de kullanırız:
\startycommalist
  \iftag{defNot}{\ycomma $\lnot$ (değilleme)}{}%
  \iftag{defAnd}{\ycomma $\land$ (ve bağlacı)}{}%
  \iftag{defOr}{\ycomma $\lor$ (veya bağlacı)}{}%
  \iftag{defIf}{\ycomma $\lif$ (!!{conditional})}{}%
  \iftag{defIff}{\ycomma $\liff$ (!!{biconditional})}{}%
  \iftag{defFalse}{\ycomma $\lfalse$ (!!{falsity})}{}%
  \iftag{defTrue}{\ycomma $\ltrue$ (!!{truth}).}{}% 
}{}

Aynı bağlaçlar çok değerli mantıklarda da kullanılır. Bununla birlikte,
örneğin ve bağlacının farklı sürümlerini aynı mantıkta içermek çoğu zaman
yararlıdır; sürümleri ayrı tutmak için de farklı simgeler gerekir. Bazı çok
değerli mantıklar klasik mantıkta karşılığı olmayan bağlaçlar da içerir.
Dolayısıyla alışıldığından biraz daha genel olacağız.

\begin{defn}
Bir \emph{önerme dili}, \emph{bağlaçlardan} oluşan bir $\Lang L$ kümesidir.
Her $\star$ bağlacının bir \emph{aritesi} vardır; aritesi $n$ olan bağlaca
\emph{$n$-konumlu} denir. Aritesi $0$ olan bağlaçlara \emph{sabitler},
aritesi $1$ olanlara \emph{tekli}, aritesi $2$ olanlara \emph{ikili} denir.
\end{defn}

\begin{ex}
Önerme mantığının standart dili $\Lang L_0$, şu bağlaçlardan (yanlarında
arite değerleriyle) oluşur: $\lfalse$~($0$), $\lnot$~($1$),
$\land$~($2$), $\lor$~($2$), $\lif$~($2$). Ele alacağımız mantıkların çoğu
bu dili kullanacaktır. Bazı mantıklar gelenek ve uzlaşım gereği bazı
bağlaçlar için farklı simgeler kullanır. Örneğin çarpım mantığında ve bağlacı
için $\land$ yerine çoğu zaman $\odot$ kullanılır. Bazen belirlenmişlik
işleci $\triangle$ (tek konumlu) gibi yeni bir işleç eklemek elverişlidir.
\end{ex}

